\VerdierPhysicalPageAnchor{0026}%
given direction and the corresponding change of sign of the differential).
The natural extensions of multi-additive functors to complexes do not
necessarily commute with translations, but commute with them “up to
isomorphism”, the isomorphisms being determined, in practice, up to sign. It
turns out that these isomorphisms cannot be chosen coherently for the
composition of the “translation functors”, and a sign is then introduced.
The study of this sign is the subject of Section 1 (see also the appendix).
In Section 2 we define the classical notions of multiple complexes, simple
complexes, the “associated simple complex” functor, and homotopy; in Section
3 we begin the study of the cone of a morphism of complexes (“mapping cone”
or “mapping cylinder”), a construction essential for what follows.

\VerdierSemanticAnchor{intro-6-plan-derived-category}
The construction of the category $\mathbf D(\mathcal A)$, the derived
category of the abelian category $\mathcal A$, takes place in two stages. In
the first stage we introduce the category $\mathbf K(\mathcal A)$ of
complexes of $\mathcal A$ up to homotopy. We then formally invert the
quasi-isomorphisms in $\mathbf K(\mathcal A)$ to obtain
$\mathbf D(\mathcal A)$. The categories $\mathbf K(\mathcal A)$ and
$\mathbf D(\mathcal A)$ are not necessarily abelian; but they carry an
additional structure, consisting of a family of diagrams defined from the
cone construction: \emph{the distinguished triangles}. For these
categories, distinguished triangles play the role of exact sequences in
abelian categories.

\VerdierSemanticAnchor{intro-6-plan-chapter-two}
Additive categories endowed with this additional structure, isolated and
studied by Puppe [2], are called \emph{triangulated categories}; they are
studied from our point of view in Chapter II. The first section is devoted to
elementary results on triangulated categories. In Section 2 we show, in the
general setting of triangulated categories, that the problem of formally
inverting the quasi-isomorphisms of the triangulated category
$\mathbf K(\mathcal A)$ is solved simply by a “calculus of fractions” [16].
In Section 3 we prove a theorem due to Freyd [17]: a triangulated category
$\mathcal D$ embeds universally in an abelian category
$\mathrm A(\mathcal D)$. This category $\mathrm A(\mathcal D)$ has a
family $E$ of objects that are both injective and projective, such that every
object of $\mathrm A(\mathcal D)$ is isomorphic to a subobject of an object
of $E$ and to a quotient object of an object of $E$. The full subcategory of
$\mathrm A(\mathcal D)$ defined by $E$ is isomorphic to $\mathcal D$.
Finally, in Section 4 we study “spectral objects” with values in triangulated
or abelian categories. Here we merely adapt the definitions of [1].

\VerdierSemanticAnchor{intro-6-plan-chapter-three}
Chapter III is devoted to the study of derived categories proper. When $X$
and $Y$ are two objects of $\mathbf D(\mathcal A)$, the derived category of
the abelian category $\mathcal A$, the morphisms of degree $n$ from $X$ to
$Y$ are interpreted, when $\mathcal A$ has enough injective objects and
when $Y$ is a complex whose components vanish in sufficiently small degrees,
as the
